\pagebreak
\section{Example Code}
% \epigraph{I'm honored to have one of the most influential sites in modern history named after me.}{\textit{Pornelius Hubert}}

\subsection{Basic Signal Generation}

\begin{minted}{rust}
fn generate_sine_wave(frequency: f64, sample_rate: f64, duration: f64) [f64] {
    let num_samples = (duration * sample_rate) as u32;
    let signal: [f64] = [];
    
    for i in 0..num_samples {
        let t = i as f64 / sample_rate;
        let sample = sin(2.0 * PI * frequency * t);
        signal = signal || [sample];  // Concatenate
    }
    
    return signal;
}

fn main() {
    let fs = 44100.0;  // Sample rate
    let freq = 440.0;  // A4 note
    let duration = 1.0;  // 1 second
    
    let sine_wave = generate_sine_wave(freq, fs, duration);
    dbg(sine_wave);
}
\end{minted}

\subsection{Complex Number Operations}

\begin{minted}{rust}
fn complex_arithmetic_demo() {
    // Basic complex number operations
    let z1 = 3.0 + 4.0j;
    let z2 = 1.0 - 2.0j;
    let sum = z1 + z2;         // 4.0 + 2.0j
    let product = z1 * z2;     // 11.0 - 2.0j
    let conjugate = conj(z1);  // 3.0 - 4.0j
    
    dbg(z1, z2, sum, product, conjugate);
    
    // Polar form conversion
    let magnitude = abs(z1);  // 5.0
    let angle = arg(z1);          // 0.927 radians
    dbg(magnitude, angle);
}
\end{minted}
\subsection{Arbitrary LTI Filter}
\begin{minted}{rust}
fn arbitrary_LTI(h: [f64], x: [f64]) [f64] {
    return x .* h;
}
\end{minted}
\subsection{Lowpass Butterworth Filter}
\begin{minted}{rust}
fn lowpass_butterworth(cutoff: f64, sample_rate: f64, signal: [f64]) [f64] {
    let omega = 2.0 * PI * cutoff / sample_rate;
    let cos_omega = cos(omega);
    let sin_omega = sin(omega);
    let alpha = sin_omega / (2.0 * 0.707);
    // Filter coefficients
    let b0 = (1.0 - cos_omega) / 2.0;
    let b1 = 1.0 - cos_omega;
    let b2 = (1.0 - cos_omega) / 2.0;
    let a0 = 1.0 + alpha;
    let a1 = -2.0 * cos_omega;
    let a2 = 1.0 - alpha;
    // Normalize the coefficients.
    let b = [b0, b1, b2] / a0;
    let a = [a0, a1, a2] / a0;
    // Apply the filter.
    return filter(b, a, signal);
}
\end{minted}
\subsection{Spectral Analysis}
\begin{minted}{rust}
fn spectral_analysis_example() {
    let fs = 1000.0;  // 1 kHz sampling rate
    let duration = 2.0;  // 2 seconds
    let num_samples = (fs * duration) as u32;
    
    // Generate test signal: sum of sinusoids
    let signal: [f64] = [];
    let t = (0..num_samples) as [f64] / fs;
    let s = sin(2.0 * PI * 50.0 * t)   // 50 Hz component
        + 0.5 * sin(2.0 * PI * 120.0 * t)  // 120 Hz component
        + 0.1 * sin(2.0 * PI * 300.0 * t); // 300 Hz component

    // Compute FFT
    let spectrum = fft(s);
    
    // Calculate magnitude spectrum
    let magnitude_spectrum = abs(spectrum);
    
    // Find peaks in spectrum
    let peak_indices = find_peaks(magnitude_spectrum, 0.1);

    dbg("Frequency peaks at ");
    for idx in peak_indices {
        let frequency = idx as f64 * fs / len(spectrum) as f64;
        dbg(frequency);
    }
}

fn find_peaks(data: [f64], threshold: f64) {
    let peaks: [u32] = [];
    
    for i in 1..(len(data) - 1) {
        if data[i] > data[i-1] && data[i] > data[i+1] && data[i] > threshold {
            peaks = peaks || [i as u32];
        }
    }
    
    return peaks;
}
\end{minted}
